% appendices/C_mi_proofs.tex
% Information-theoretic proofs for Chapter 6 (MI hallucination detection).
% Ported from the HalluLens source paper's Appendix A ("Theorems to be added"):
% ../sources/hallulens_mi_hallucination/03_method.md (Theorems section).
% The source's A.* labels are RENUMBERED to this appendix to avoid colliding
% with Appendix A (Label Blindness): the source's Lemma A.2 / Lemma A.3 /
% Proposition A.1 / Corollary A.1.1 / Corollary A.2.1 become the lemmas,
% proposition, and corollaries of this appendix. The Proposition STATEMENT lives
% in the chapter body (\Cref{prop:hd:supcon-mi}); this appendix supplies its
% proof, the two supporting lemmas, and the two corollaries.
%
% NOTE on the number gate (make lint): lint.py strips math line-by-line and does
% not recognize multi-line display environments, so every display equation here
% is kept on a single physical line inside a stripped environment.
\chapter{Information-Theoretic Proofs for Hallucination Detection}
\label{app:mi-hallucination-proofs}
% Short running head: the full title overflows the page header by ~17pt. Full
% title kept in the ToC.
\chaptermark{Proofs for Hallucination Detection}

This appendix supplies the formal counterparts of the feasibility argument made
in the chapter on information-theoretic hallucination detection
(\Cref{sec:hd:method}). The load-bearing result is
\Cref{prop:hd:supcon-mi}, whose statement appears in the chapter body and whose
proof is given here; it rests on two standard lemmas---a data-processing chain
(\Cref{lem:hd:dpi}) and the InfoNCE variational lower bound
(\Cref{lem:hd:infonce})---which we restate in the notation of the chapter so the
proposition has named results to compose. We also record two corollaries: one
strengthening the cross-layer information claim to the response variable
(\Cref{cor:hd:coinfo}), and one that turns the proposition's asymmetry into a
falsifiable prediction about symmetric supervised contrastive learning
(\Cref{cor:hd:symcon}), confirmed empirically in \Cref{sec:hd:ablations}.

Throughout we use the notation of \Cref{sec:hd:method}. The frozen language model
$M$ has $L$ residual-stream layers; for an input--generation pair $(x,y)$ we
write $h_\ell = h_\ell(x,y) \in \mathbb{R}^{T_y \times d}$ for the full
per-token activation sequence at layer $\ell$, and $z_\ell = f_\theta(h_\ell)$
for its compressed embedding. A \emph{view} of a generation is the embedding of
one layer drawn from the extraction band; two views of the same generation are
$z_a = f_\theta(h_{\ell_a})$ and $z_b = f_\theta(h_{\ell_b})$ for a sampled layer
pair $\ell_a < \ell_b$. The operational label is
$y_{\mathrm{label}} \in \{0,1\}$, with $y_{\mathrm{label}}=0$ for a truthful
generation and $y_{\mathrm{label}}=1$ for a hallucinated one.

\section{Data-Processing Chain for the Residual Stream}
\label{sec:hd-proofs:dpi}

The residual stream is deterministic and additive: each layer's activation is
its predecessor plus the output of the intervening block. This gives both the
positivity of cross-layer mutual information that the feasibility argument starts
from, and the direction of the data-processing inequality that lets a positive
mutual information measured at the \emph{embedding} level certify a positive
mutual information at the \emph{activation} level.

\begin{lemma}[Cross-layer data-processing chain]\label{lem:hd:dpi}
Fix a prompt and let $\ell_a < \ell_b$ index two residual-stream layers of $M$.
\emph{(i)~Forward direction.} Because $h_{\ell_b}$ is a deterministic function of
$h_{\ell_a}$ and the intervening residual blocks, with the residual connections
carrying $h_{\ell_a}$ forward additively, the Markov-free determinism gives
$I(h_{\ell_a}; h_{\ell_b}) = H(h_{\ell_b}) > 0$ for any input distribution under
which $h_{\ell_b}$ is non-degenerate. \emph{(ii)~Backward direction.} For any
(possibly stochastic) encoder $f_\theta$ applied layerwise, the embeddings
$z_a = f_\theta(h_{\ell_a})$ and $z_b = f_\theta(h_{\ell_b})$ satisfy
\begin{equation} I(z_a; z_b) \;\le\; I(h_{\ell_a}; h_{\ell_b}). \label{eq:hd:dpi} \end{equation}
\end{lemma}

\begin{proof}
\emph{(i)} Determinism of the forward pass gives
$H(h_{\ell_b} \mid h_{\ell_a}) = 0$, so
$I(h_{\ell_a}; h_{\ell_b}) = H(h_{\ell_b}) - H(h_{\ell_b}\mid h_{\ell_a}) = H(h_{\ell_b})$,
which is strictly positive whenever $h_{\ell_b}$ is not almost surely constant.
\emph{(ii)} The variables form the Markov chain
$z_a \leftarrow h_{\ell_a} \rightarrow h_{\ell_b} \rightarrow z_b$, since $z_a$
depends on $(h_{\ell_a},h_{\ell_b})$ only through $h_{\ell_a}$ and $z_b$ only
through $h_{\ell_b}$. The data-processing inequality applied to this chain yields
$I(z_a; z_b) \le I(h_{\ell_a}; h_{\ell_b})$.
\end{proof}

\begin{remark}
The two directions play complementary roles in the chapter. Part~(i) supplies the
\emph{architectural} guarantee that the joint $(h_{\ell_a},h_{\ell_b})$ carries
shared information at all. Part~(ii) is what makes trainability evidential: a
positive mutual information observed between the learned embeddings is a lower
bound witness for the activation-level quantity, never an artifact of the
encoder---the inequality~\eqref{eq:hd:dpi} runs in the favourable direction.
\end{remark}

\section{InfoNCE Variational Lower Bound}
\label{sec:hd-proofs:infonce}

\begin{lemma}[InfoNCE lower bound on mutual information]\label{lem:hd:infonce}
Let $(\mathrm{view}_a, \mathrm{view}_b)$ be paired views with $K-1$ independent
negatives, and let $\mathcal{L}_{\mathrm{NCE}}$ be the InfoNCE loss evaluated
with a critic over those $K$ candidates. Then
\begin{equation} \mathcal{L}_{\mathrm{NCE}} \;\ge\; \log K - I(\mathrm{view}_a; \mathrm{view}_b), \qquad\text{equivalently}\qquad I(\mathrm{view}_a; \mathrm{view}_b) \;\ge\; \log K - \mathcal{L}_{\mathrm{NCE}}. \label{eq:hd:infonce} \end{equation}
In particular, if $I(\mathrm{view}_a;\mathrm{view}_b) = 0$ then
$\mathcal{L}_{\mathrm{NCE}} \ge \log K$, attained by uniform negatives.
\end{lemma}

\begin{proof}
This is the standard contrastive lower bound of~\citet{oord2018representation};
we use the tightened analysis of~\citet{poole2019variational}. Writing the
InfoNCE objective as the expected log-softmax of the critic over one positive and
$K-1$ negatives, Jensen's inequality on the partition function bounds the
expected score gap by $I(\mathrm{view}_a;\mathrm{view}_b) - \log K$ from above,
which rearranges to~\eqref{eq:hd:infonce}. The $\log K$ floor follows because the
optimal critic under zero mutual information cannot distinguish the positive from
the $K-1$ negatives, so every candidate receives equal posterior mass $1/K$ and
the loss equals $\log K$.
\end{proof}

The contrapositive of the floor is the chapter's feasibility punchline: a loss
that descends below $\log K$ certifies $I(\mathrm{view}_a;\mathrm{view}_b) > 0$,
and by \Cref{lem:hd:dpi}(ii) this transfers to the activation level. The
existence of a trainable probe is therefore the empirical content of the
feasibility claim, not merely consistent with it.

\section{Cross-Layer/Response Co-Information}
\label{sec:hd-proofs:coinfo}

\Cref{lem:hd:dpi}(i) states only that the cross-layer joint carries shared
information. The following corollary ties that shared information to the response
variable $y$ as well, so the joint is non-orthogonal to what the model generated.

\begin{corollary}[Cross-layer/response co-information under downward determinism]\label{cor:hd:coinfo}
Suppose $h_{\ell_b}$ is a deterministic function of $h_{\ell_a}$ given the prompt
(residual-stream determinism). Then the three-way interaction information
satisfies
\begin{equation} I(h_{\ell_a}; h_{\ell_b}; y) \;:=\; I(h_{\ell_a}; h_{\ell_b}) - I(h_{\ell_a}; h_{\ell_b} \mid y) \;=\; I(h_{\ell_b}; y), \label{eq:hd:coinfo} \end{equation}
which is strictly positive whenever $h_{\ell_b}$ is a non-constant function of
$y$.
\end{corollary}

\begin{proof}
Downward determinism gives
$H(h_{\ell_b}\mid h_{\ell_a}) = H(h_{\ell_b}\mid h_{\ell_a}, y) = 0$, hence
$I(h_{\ell_a}; h_{\ell_b}) = H(h_{\ell_b})$ and
$I(h_{\ell_a}; h_{\ell_b}\mid y) = H(h_{\ell_b}\mid y)$. Substituting both into
the definition of the interaction information gives
$I(h_{\ell_a}; h_{\ell_b}; y) = H(h_{\ell_b}) - H(h_{\ell_b}\mid y) = I(h_{\ell_b}; y)$,
which is positive precisely when $h_{\ell_b}$ depends non-trivially on $y$.
\end{proof}

\begin{remark}[Surviving caveat]
Under the operational definition
$y_{\mathrm{label}} = g_{\mathrm{eval}}(x,y)$---deterministic in $(x,y)$ for a
fixed evaluator---\Cref{cor:hd:coinfo} bounds $I(h_{\ell_a};h_{\ell_b};y)$ in
terms of the \emph{response} variable $y$, and thereby gives only an
\emph{upper} bound on $I(h_{\ell_a};h_{\ell_b};y_{\mathrm{label}})$, not a lower
bound. Whether the label-relevant features survive the lossy compression that a
single layer applies to $(x,y)$ is an empirical question, settled by the
trainability punchline rather than by this corollary.
\end{remark}

\section{Asymmetric Supervised Contrastive Bound}
\label{sec:hd-proofs:supcon}

We now prove the proposition stated in the chapter body
(\Cref{prop:hd:supcon-mi}): the asymmetric \texttt{ignore\_label} scheme makes
the optimized contrastive bound \emph{one-sided}, bounding the mutual information
between the learned features and the truthful class only, while the
hallucinated-anchor term reverts to instance-level view consistency. The proof
turns on a decomposition of the loss by anchor class.

Let $Z = f_\theta(h)$ denote the learned feature and recall the asymmetric
scheme: for a truthful anchor ($y_{\mathrm{label}}=0$) the positives are the
other truthful samples in the batch together with the anchor's own other-layer
views, and all hallucinated samples are negatives; for a hallucinated anchor
($y_{\mathrm{label}}=1$) the positives are \emph{only} the anchor's own
other-layer views, and everything else in the batch---including other
hallucinated samples---is a negative.

\begin{proof}[Proof of \Cref{prop:hd:supcon-mi}]
\emph{Step~$1$ (decomposition by anchor class).}
Because each anchor contributes an independent additive term to the supervised
contrastive objective, the total loss splits along the anchor label,
\begin{equation} \mathcal{L}_{\mathrm{SupCon\text{-}asym}} \;=\; \mathcal{L}_{\mathrm{truth}} + \mathcal{L}_{\mathrm{hall}}, \label{eq:hd:loss-split} \end{equation}
where $\mathcal{L}_{\mathrm{truth}}$ collects the truthful anchors (positives:
other truthful samples and own other-layer views) and
$\mathcal{L}_{\mathrm{hall}}$ collects the hallucinated anchors (positives: own
other-layer views only).

\emph{Step~$2$ (truthful anchors bound a class-conditioned mutual information).}
On the truthful subset, grouping the positives by the label
$\mathbb{1}[y_{\mathrm{label}}=0]$ is exactly the supervised-contrastive
specialization of~\citet{khosla2020supervised}: the InfoNCE numerator's
instance-level similarity is replaced by a class-averaged similarity over the
truthful cluster. Composing this class grouping with \Cref{lem:hd:infonce} and
the data-processing chain on the class variable tightens the truthful-anchor term
into a lower bound on the mutual information between the features and the truthful
indicator,
\begin{equation} -\,\mathcal{L}_{\mathrm{truth}} \;\le\; I\bigl(Z; \mathbb{1}[y_{\mathrm{label}}=0]\bigr) - \log K + c, \label{eq:hd:truth-bound} \end{equation}
with $c$ the constant absorbing the class-averaging normalization. The features
are pushed to distinguish ``in the truthful cluster'' from ``everything else.''

\emph{Step~$3$ (hallucinated anchors revert to instance-level InfoNCE).}
For a hallucinated anchor the positive set is its own other-layer views only, so
no class grouping occurs: the numerator stays instance-level and
\Cref{lem:hd:infonce} bounds $\mathcal{L}_{\mathrm{hall}}$ by the mutual
information between paired views of an individual hallucinated generation---view
consistency---and contributes \emph{no} term in
$I(Z; \mathbb{1}[y_{\mathrm{label}}=1])$, because the hallucinated samples are
never grouped as mutual positives.

\emph{Step~$4$ (sum).}
Adding the two terms via~\eqref{eq:hd:loss-split}, the optimized objective lower
bounds $I(Z; \mathbb{1}[y_{\mathrm{label}}=0])$ from Step~$2$ plus an
instance-level view-consistency term from Step~$3$, with no class-mutual-
information term on the hallucinated side. The bound is therefore one-sided: it
constrains the truthful class only. This is the claim.
\end{proof}

\begin{remark}[What the proposition does not establish]
Three honest negatives accompany the bound. First, it is not claimed tight; the
looseness of InfoNCE bounds is analyzed by~\citet{poole2019variational}. Second,
it does not imply sufficiency for hallucination detection---that the truthful
cluster is separable from hallucinations---which is settled empirically by the
trainability punchline of \Cref{sec:hd:method}. Third, the role of the
hallucinated-anchor term as a regularizer that prevents collapse of the
hallucinated features is an empirical observation, not a consequence of the
bound.
\end{remark}

\section{Symmetric Supervised Contrast Is Contraindicated}
\label{sec:hd-proofs:symcon}

The asymmetry in \Cref{prop:hd:supcon-mi} is not an arbitrary implementation
choice. The following corollary shows that the symmetric scheme---labelling both
classes as mutual-positive clusters---injects a gradient that actively opposes
the truthful-cluster signal when the hallucinated class lacks coherent latent
structure, and predicts the empirical degradation reported in
\Cref{sec:hd:ablations}.

\begin{corollary}[Symmetric SupCon contraindication]\label{cor:hd:symcon}
Suppose the hallucinated positive set under \emph{symmetric} supervised contrast
has near-uniform pairwise similarity, i.e. its intra-class InfoNCE numerator
terms are within $o(1)$ of the inter-class terms. Then the expected
hallucinated-class gradient approaches zero in magnitude but retains a residual
component directed \emph{away} from the truthful cluster, because the InfoNCE
denominator for hallucinated anchors includes the truthful negatives. The net
effect is a gradient contribution negatively correlated with the truthful-anchor
gradient, with magnitude bounded by the hallucinated-class loss mass times the
class-imbalance ratio. Consequently the symmetric scheme does not reduce to the
asymmetric scheme of \Cref{prop:hd:supcon-mi} under any choice of temperature or
loss weight: the two differ in the construction of the hallucinated positive set,
not by a scale factor.
\end{corollary}

\begin{proof}
Reusing the decomposition~\eqref{eq:hd:loss-split}, the only change under the
symmetric scheme is that hallucinated anchors acquire the other hallucinated
samples as positives. Write the hallucinated-anchor gradient as the difference of
a numerator pull toward the hallucinated positives and a denominator push away
from all negatives. Under the near-uniform-similarity hypothesis the numerator
pull averages to a vector of vanishing magnitude (all hallucinated pairwise
similarities coincide to first order, so the softmax-weighted positive mean
cancels against the anchor). The denominator term does not vanish: it retains the
truthful negatives, producing a residual component pointing away from the
truthful cluster. Projected onto the truthful-anchor gradient
of~\eqref{eq:hd:truth-bound}, this residual has negative expected inner product,
and its norm is controlled by the hallucinated-class loss mass scaled by the
ratio of hallucinated to truthful batch counts. Since the construction of the
hallucinated positive set---not a multiplicative constant---is what changes, no
temperature or weight rescaling recovers the asymmetric objective.
\end{proof}

\begin{remark}[Confirmed prediction]
\Cref{cor:hd:symcon} makes a falsifiable prediction: relabelling the hallucinated
side as a mutual-positive cluster should \emph{degrade} detection rather than
leave it unchanged. The Variant~$4$ ablation of \Cref{sec:hd:ablations}, which
supervises both classes symmetrically, confirms it---an AUROC drop of roughly ten
points relative to the one-class asymmetric scheme, on both supported models.
Theory predicted the failure mode; the ablation observed it.
\end{remark}
